\section{Research Projects}

% \underline{NASA DEVELOP Projects}:
\outerlist{

\entrybigwrap
    {{\textbf{Measuring Crop Yields in Northern Brazil During El Niño Conditions to Evaluate Trends in Agricultural Production and Support Crop Forecasting}}{\hfill NOAA NCEI}}
    {{\textsc{Role:} Research Fellow}{\hfill Jan.\textendash March 2024}}
    %\innerlist{
        %\entry{\underline{\textsc{Advisors:}} Dr. Jesse \textsc{Bell}, Dr. Jesse \textsc{Berman}, Dr. Babak \textsc{Fard}, Ronald \textsc{Leeper}, Jared \textsc{Rennie}, Molly \textsc{Woloszyn}}
        %\entry{\underline{\textsc{Project Team:}} Greta Bolinger, Tallis Monteiro, Abby Sgan, Cristina Villalobos-Heredia, Taylor West}
        % \entry{Collaborated with the USDA Foreign Agriculture Service International Production Assessment Division and World Agricultural Outlook Board to evaluate changes in crop production in northern Brazil with El Niño-Southern Oscillation using NDVI, precipitation, and temperature data.}
%}

\entrybigwrap
    {{\textbf{Multi-Hazard Approach to Mapping Flood Susceptibility and Vulnerability in Kentucky}}}
    {{NOAA NCEI}{\hfill Jan.\textendash March 2024}}
    % {{\textsc{Role:} Research Fellow}{\hfill Jan\textendash March. 2024}}
    % \innerlist{
        %\entry{\underline{\textsc{Advisors:}} Dr. Jesse \textsc{Bell}, Dr. Jesse \textsc{Berman}, Dr. Babak \textsc{Fard}, Ronald \textsc{Leeper}, Jared \textsc{Rennie}, Molly \textsc{Woloszyn}}
        %\entry{\underline{\textsc{Project Team:}} Greta Bolinger, Tallis Monteiro, Abby Sgan, Cristina Villalobos-Heredia, Taylor West}
        % \entry{Modeled flood risk in Kentucky in partnership with the Kentucky Climate Center and NOAA National Weather Service Jackson and Paducah Forecast Offices based on flood and vulnerability susceptibility factors. Mapped antecedent soil moisture conditions prior to flood events in Kentucky in 2022 and 2023.}
% }

\entrybigwrap
    {{\textbf{Methods for Monitoring Rhode Island Habitats: Contributing to a Framework for Targeted Conservation and Management}}}
    {{NOAA NCEI}{\hfill Jan.\textendash March 2024}}
    % {{\textsc{Role:} Research Fellow}{\hfill Jan\textendash March. 2024}}
    % \innerlist{
        %\entry{\underline{\textsc{Advisors:}} Dr. Jesse \textsc{Bell}, Dr. Jesse \textsc{Berman}, Dr. Babak \textsc{Fard}, Ronald \textsc{Leeper}, Jared \textsc{Rennie}, Molly \textsc{Woloszyn}}
        %\entry{\underline{\textsc{Project Team:}} Greta Bolinger, Tallis Monteiro, Abby Sgan, Cristina Villalobos-Heredia, Taylor West}
        % \entry{Supervised four-person team that mapped historical land cover cover change and modeled future land cover change in Rhode Island to inform land conservation and management efforts by the Audubon Society of Rhode Island.}
% }

\entrybigwrap
    {{\textbf{Monitoring Trends in Air Quality During a Drought Case Study to Improve Public Health Response to Drought Threats}}}
    {{NOAA NCEI}{\hfill June\textendash Aug. 2023}}
    % {{\textsc{Role:} Research Fellow}{\hfill June\textendash Aug. 2023}}
    % \innerlist{
        % \entry{\underline{\textsc{Advisors:}} Dr. Jesse \textsc{Bell}, Dr. Jesse \textsc{Berman}, Dr. Babak \textsc{Fard}, Ronald \textsc{Leeper}, Jared \textsc{Rennie}, Molly \textsc{Woloszyn}}
        %\entry{\underline{\textsc{Project Team:}} Greta Bolinger, Tallis Monteiro, Abby Sgan, Cristina Villalobos-Heredia, Taylor West}
        % \entryextra{Proposed project and supervised five-person team to analyze trends in PM2.5, PM10, and aerosol optical depth during drought evolution in partnership with local health departments in Oregon and Washington.}
        %\entryextra{\underline{\textsc{Publication:}}} 
% }

\entrybigwrap
    {{\textbf{Using Earth Observations to Evaluate Snow Variability through a Climatological Analysis to Support Ecological Monitoring in Northeast Alaska}}}
    {{NOAA NCEI}{\hfill Jan.\textendash March 2023}}
    % {{\textsc{Role:} Research Fellow}{\hfill Jan.\textendash March 2023}}
    % \innerlist{
        % \entry{\underline{\textsc{Advisors:}} Dr. Bob \textsc{Bolton}, Dr. Jessica \textsc{Cherry}, Julian \textsc{Dann}, Molly \textsc{Woloszyn}}
        %\entry{\underline{\textsc{Project Team:}} Kristin Anderson, Omeed Arooji, Isabella Chittumuri, Thomas Germann}
        % \entryextra{Proposed study with U.S. Fish and Wildlife Service, Alaska Climate Adaptation Science Center, and University of Alaska Fairbanks and supervised four-person team that assessed snow cover variability to support conservation in the Arctic National Wildlife Refuge and National Petroleum Reserve.}
% }

\entrybigwrap
    {{\textbf{Evaluating the Role of Soil Moisture in Determining Vegetation Health, Fuel Loads, and Wildfires in the Gatlinburg and Beatty Wildfires}}}
    {{NOAA NCEI}{\hfill Sept.\textendash Nov. 2022}}
    % {{\textsc{Role:} Research Fellow}{\hfill Sept.\textendash Nov. 2022}}
    % \innerlist{
        % \entry{\underline{\textsc{Advisors:}} Ronald \textsc{Leeper} and Molly \textsc{Woloszyn}}
        %\entry{\underline{\textsc{Project Team:}} William Hadley, Daniel Littleton, Kelli Roberts}
        % \entryextra{Supervised three-person team that compared antecedent soil moisture, vegetation health, and fire fuel indices preceding fire events to assess utility in fire potential monitoring.}
        %\entryextra{\underline{\textsc{Publication:}} Evaluating the role of soil moisture in determining vegetation health, fuel loads, and wildfires in the Gatlinburg and Beatty Wildfires}
% }

\entrybigwrap
    {{\textbf{Monitoring Trends in Tree Defoliation Due to \textit{Lymantria dispar} Outbreaks to Predict Future Hardwood Tree Mortality and Health Impacts}}}
    {{NASA Ames Research Center}{\hfill Jan.\textendash April 2022}}
    % {{\textsc{Role:} Researcher}{\hfill Jan.\textendash April 2022}}
    %Monitoring Trends in Tree Defoliation Due to Lymantria dispar Outbreaks to Predict Future Hardwood Tree Mortality and Health Impacts
    % \innerlist{
        % \entry{\underline{\textsc{Advisor:}} Dr. Juan \textsc{Torres-Perez}}
        %\entry{\underline{\textsc{Co-Investigators:}} Morgan Dean, Jacob Orser, Seamore Zhu}
        % \entryextra{Modeled tree defoliation due to the invasive month \textit{Lymantria dispar} using Landsat 7/8, Sentinel-2, SMAP, and Terra MODIS satellite data to support forest management.}
% }

\entrybigwrap
    {{\textbf{Utilizing NASA Earth Observations to Map Suitable Venus Flytrap Habitat in an Effort to Inform Conservation, Seed Banking, and Reintroduction in the Carolina Coastal Plain and Sandhills Regions}}}
    {{UGA Center for Geospatial Research}{\hfill Sept.\textendash Nov. 2021}}
    % {{\textsc{Role:} Researcher}{\hfill Sept.\textendash Nov. 2021}}
    % \innerlist{
        % \entry{\underline{\textsc{Advisor:}} Dr. Marguerite \textsc{Madden}}
        %\entry{\textsc{\underline{Co-Investigators:}} Jayne Lampley, Monika Rock, Ashna Siddhi, Seamore Zhu}
        % \entryextra{Modeled present-day and future suitable \textit{Dionaea muscipula} habitat to inform seed banking and reintroduction efforts by the North Carolina Botanical Garden and North Carolina Natural Heritage Program.}
% }

\entrybigwrap
    {{\textbf{Assessing the effect of eastern hemlock (\textit{Tsuga canadensis}) {decline from hemlock woolly adelgid} (\textit{Adelges tsugae}) {infestation on ectomycorrhizal colonization and growth of red oak} (\textit{Quercus rubra}) {seedlings}}}{\hfill Depts. of Biology, UNC Asheville and Warren Wilson College}}
    {{\textsc{Role:} Undergraduate Researcher}{\hfill May 2019\textendash May 2020}}
    % \innerlist{
        % \entry{\textsc{\underline{Advisors:}} Dr. Jonathan L. \textsc{Horton} and Dr. Alisa A. \textsc{Hove}}
        % \entryextra{Investigated the effect of \textit{Tsuga canadensis} decline from \textit{Adelges tsugae} infestation by comparing mycorrhizal community assemblages and seedling growth between a healthy and declining \textit{T. canadensis} stand. Measured seedling growth, biomass, and ectomycorrhizal colonization of seedling roots, collected ectomycorrhizal morphotypes, and characterized ectomycorrhizal taxonomic richness using DNA barcoding analyses.}
    %Eastern hemlock (Tsuga canadensis (L.) Carrie`re) is a foundation species in eastern North American forests, providing critical habitats for a number of species. These trees are experiencing widespread decline due to the spread of hemlock woolly adelgid (HWA: Adelges tsugae Annand Order Hemiptera) into their range, potentially resulting in the disappearance of hemlocks from eastern forests. Hemlock dieback can lead to cascading effects on associated ecosystems, including belowground, mycorrhizal fungal communities. Ectomycorrhizal fungi (EM), which are mutualistic with many tree species and provide nutrients to plant hosts, are known to colonize hemlock as well as neighboring tree species at lower levels following HWA infection. This study investigated the effect of hemlock decline from HWA infestation on mycorrhizal communities, as inferred from colonization on northern red oak (Quercus rubra L.) ‘‘bait’’ seedlings grown near ‘‘host’’ hemlock trees. Hemlock health surveys were conducted in healthy (Carl Sandburg Home National Historic Site – CARL) and declining (Warren Wilson College – WWC) stands in western North Carolina, and host trees were paired between stands based on diameter. In each stand, northern red oak seedlings were planted within a meter of host hemlocks in early summer and allowed to grow for 8 w, when they were harvested. Seedling growth and dry biomass were recorded at harvest and roots were sampled for mycorrhizal colonization frequencies. Different mycorrhizal morphotypes were collected from seedling roots for subsequent DNA barcoding analyses to characterize EM taxonomic richness to compare mycorrhizal community assemblages between the two stands. Mycorrhizal colonization frequencies (percentage of the total number of EM-colonized root tips per seedling) and growth in seedling height were significantly greater at CARL than WWC, suggesting healthy hemlock stands are more favorable for oak seedling growth than declining stands. Moreover, a greater proportion of seedlings grown in the healthy stand were colonized by EM, indicating EM assemblages differ between a healthy and a declining hemlock stand. Differences between EM communities corresponded with altered seedling growth allocation, as seedlings in the declining stand had higher root to shoot ratios with reduced stem height, but showed greater investment in root biomass and stem diameter growth. We conclude EM communities differ between a healthy and declining hemlock stands, and changes in EM communities following hemlock dieback may affect the growth of replacement species.
% }

\entrybigwrap
    {{\textbf{Factors contributing to reproductive success and failure in Virginia spiraea (\textit{Spiraea virginiana}), a federally threatened shrub}}}
    {{Dept. of Biology, UNC Asheville}{\hfill May 2019\textendash May 2020}}
    % {{\textsc{Role:} Research Fellow}{\hfill May 2019\textendash May 2020}}
    % \innerlist{
        % \entry{\textsc{\underline{Advisor:}} Dr. Jennifer Rhode \textsc{Ward}}
        % \entryextra{Informed conservation and restoration of \textit{S. virginiana} by comparing reproductive effort and output among populations, manipulating breeding systems to assess barriers to seed production, identifying insect visitors to ascertain potential pollinators, measuring physiological response to light availability, and testing hybridization potential with Japanese spiraea (\textit{Spiraea japonica}), an invasive congener.}
    %Virginia spiraea (Spiraea virginiana) is a perennial shrub with both sexual and asexual modes of reproduction, found in riparian habitats throughout the Appalachians. The species was listed as federally threatened in 1990 due to hydrologic disturbance, and populations might suffer from inbreeding depression. In this series of experiments, we examined factors that could challenge the species’ survival. We compared reproductive effort and output among populations, manipulated breeding systems to see their effects on seed production, identified insect visitors to ascertain potential pollinators, used physiological measurements in experimental gardens to determine whether the species is shade-tolerant, and tested the potential for hybridization with Japanese spiraea (Spiraea japonica), an invasive congener. Results showed variation in reproductive effort and seed production among genetic individuals and populations, and analyses of breeding systems are ongoing. Guilds of insect visitors varied widely in their potential pollination roles; because spiraea pollen loads were negligible in ants and flies, it is likely that other groups (beetles, bees, wasps) are spiraea’s major pollinators. Virginia spiraea’s photosynthetic rate did not saturate until full light, and unshaded plants produced more corymbs, suggesting that removal of woody species could increase spiraea’s performance under field conditions. Finally, although reproductive effort was high in Japanese spiraea, cross-pollinations between native and invasive spiraea produced no seeds. Results will inform efforts to restore and augment populations of Virginia spiraea throughout its range, and will ensure the continued persistence of existing clones.
    %This project is a collaboration between the University of North Carolina at Asheville and the Botanical Gardens at Asheville to research Spiraea virginiana (Virginia spiraea) reproduction, breeding, and physiological responses to light availability and vegetatively propagate it for outplanting at the Asheville Botanical Gardens. The work is also in collaboration with the regional office of the USFWS. The Botanical Gardens at Asheville had a clone of a S. virginiana plant that was originally collected in Buncombe County, NC and maintained by the Arnold Arboretum. Unfortunately, that clone died during the winter of 2018-2019, and no genetically related individuals are extant in the wild. Due to the demise of its clone, the Botanical Gardens at Asheville requests a propagule of plant AA# 774-88*A, which was originally collected in Buncombe County, for outplanting at the Asheville Botanical Gardens. This will serve preservation and educational purposes.
% }

\entrybigwrap
    {{\textbf{A CEREUS (Consortium Exchanging Research Experiences for Undergraduate Students) Approach to Assessing Ecological Responses in the Southern Appalachians to Environmental Change}}}
    {{Dept. of Biology, UNC Asheville}{\hfill Jan. 2019\textendash May 2020}}
    % {{\textsc{Advisor:} Dr. Jennifer Rhode \textsc{Ward}, Dr. H. David \textsc{Clarke}, & Dr. Jonathan \textsc{Horton}}{\hfill Jan. 2019\textendash May 2020}}
    %{{\textsc{Role:} Research Assistant}{\hfill Jan. 2019\textendash May 2020}}
    % \innerlist{
       % \entry{\textsc{\underline{Principal Investigator:}} Dr. Jennifer Rhode \textsc{Ward}}
        % \entry{\textsc{\underline{Co-Principal Investigators:}} Dr. H. David \textsc{Clarke} and Dr. Jonathan \textsc{Horton}}
        % \entryextra{Investigated environmental change by monitoring Southern Appalachian plant phenology, contributing observations to the National Phenology Network Nature's Notebook Citizen Science database, and studying effects of invasive exotic removal techniques on regeneration of native and nonnative plants.}
    %Phenology, the study of seasonal biological events, tracks the effects of both year to year climate variability and long-term environmental change on both local and global ecology. It has considerable potential for serving as a base for research studies by undergraduate classes and developing cross-campus studies. Several recent studies, including the President's Council of Advisors on Science and Technology (2014), and Vision and Change in Undergraduate Biology Education (a 2011 document representing the views of over 500 biology faculty and administrators) cite the importance of undergraduate research for retaining science majors and helping students better understand key concepts and develop crucial competencies. Four institutions, the University of North Carolina at Asheville, Appalachian State University, East Tennessee State University, and Warren Wilson College, are combining forces to create common gardens and to develop modules that can guide students' research studies to address potential impacts of global change at multiple levels of the biological hierarchy, from genes to ecosystems. The cooperating institutions are unique in character, diverse in their student bodies, and located in varying climate zones. The project will contribute to our knowledge about ecological issues and will expand the ways in which institutions can combine forces to increase the STEM knowledge and competencies of their students. Environmental change issues to be investigated include monitoring effects of invasive exotic removal techniques on herbs, shrubs, vines, and trees within the areas studied (studies for lower division courses), and exploring questions about the spread of potentially invasive exotic plants, using a combination of GIS (Global Information Systems) mapping, herbarium collections, and online databases (upper division courses). 
    %The collaborating institutions will create and beta-test modules, identify and mitigate barriers to successful implementation of curricular changes, and disseminate curricular changes. The modules will have a direct impact on the learning of over 2000 undergraduates per year, including non-STEM majors. They will be developed and partially administered by undergraduate and graduate research students, and will be instrumental in developing pedagogical expertise for faculty members. The ability of modules to impart botanical knowledge while encouraging higher-order cognitive processes, advancing quantitative literacy, teaching analytical techniques, honing scientific communication skills, cultivating positive student attitudes towards plants and STEM, and improving persistence and graduation in STEM majors will be assessed using a mixed method approach (quantitative and qualitative data collection). Results will begin to document the effects of global change on an understudied but important bioregion while developing a future workforce with solid STEM training.
    %This project is funded jointly by the Directorate for Biological Sciences, Division of Biological Infrastructure and the Directorate for Education and Human Resources, Division of Undergraduate Education, in support of efforts to address the challenges posed in Vision and Change in Undergraduate Education: A Call to Action http://visionandchange.org/finalreport/.
% }

%\entrybigwrap
%{{\textbf{Southern Appalachian Cooperative Ecosystems Studies Unit, U.S. Park Service}}{\hfill Carl Sandburg National Historic Site}}
    %{{\textsc{Role:} Research Assistant}{\hfill May 2019.\textendash May 2020}}
    %\innerlist{
        %\entry{\textsc{\underline{Principal Investigator:}} Dr. Jonathan L. \textsc{Horton}}
        %\entryextra{Assessed efficacy of biological and chemical for hemlock woolly adelgid infestation by monitoring hemlock tree canopy health, hemlock woolly adelgid infestation rates, and \textit{Laricobius} spp. predator beetle densities.}
%}
}